% Espaçamento de linha
\linespread{1.5}

% espaçamento dos resumos
\setlength{\absparsep}{18pt}

% Margens
\setlrmarginsandblock{3cm}{2cm}{*}
\setulmarginsandblock{3cm}{2cm}{*}
\checkandfixthelayout

% Tabelas
\setlength{\extrarowheight}{0pt}
\renewcommand{\arraystretch}{0.9}

% Parágrafos
\setlength{\parindent}{1.3cm}     % O tamanho do parágrafo
\setlength{\parskip}{0.2cm}       % Controle do espaçamento entre um parágrafo e outro

% Controle de linhas orfãs ou viúvas:
% \clubpenalty=10000
% \widowpenalty=10000

% Fonte de capítulos
% \renewcommand{\ABNTEXchapterfont}{%
%     \bfseries\fontspec{Calibri}\MakeUppercase
% }

% Redefinição do formato de capítulo
\makeatletter
\renewcommand{\ABNTEXchapterfont}{\bfseries\fontsize{14pt}{24pt}\selectfont} % negrito, 14pt
\renewcommand{\ABNTEXchapterfontsize}{\fontsize{14pt}{24pt}\selectfont}

% remove número automático antes do título
\renewcommand{\printchapternum}{}

% centraliza e define formato
\renewcommand{\printchaptertitle}[1]{%
  \begin{center}
    \ABNTEXchapterfont
    \ifnum\value{chapter}>0
      \MakeUppercase{CAPÍTULO \thechapter: #1}%
    \else
      \bfseries\fontsize{16pt}{24pt}\selectfont\MakeUppercase{#1}%
    \fi
  \end{center}%
}
\makeatother


\definecolor{meu_rosa}{HTML}{FF69B4} % HotPink



% \renewcommand{\ABNTEXchapterfontsize}{\Large}

% ----------------------------------------------------------
% Configurações de cores para o código R
% ----------------------------------------------------------

\lstset{ %
  language=R,                                       % the language of the code
  backgroundcolor=\color{white},                    % choose the background color. You must add \usepackage{color}
  basicstyle=\footnotesize\ttfamily,                             % the size of the fonts that are used for the code
  breakatwhitespace=false,                          % sets if automatic breaks should only happen at whitespace
  breaklines=true,                                  % sets automatic line breaking
  captionpos=t,                                     % sets the caption-position to Top
  commentstyle=\color{dkgreen},                     % comment style
  deletekeywords={...},                             % if you want to delete keywords from the given language
  escapeinside={\%*}{*)},                           % if you want to add a comment within your code
  frame=single,                                     % adds a frame around the code
  keywordstyle=\color{blue},                        % keyword style
  morekeywords={*,...},                             % if you want to add more keywords to the set
  numbers=left,                                     % where to put the line-numbers
  numbersep=5pt,                                    % how far the line-numbers are from the code
  numberstyle=\tiny\color{mygray},                    % the style that is used for the line-numbers
  postbreak=\mbox{\textcolor{red}{$\hookrightarrow$}\space},
  rulecolor=\color{black},                          % if not set, the frame-color may be changed on line-breaks within not-black text (e.g. commens (green here))
  showspaces=false,                                 % show spaces adding particular underscores
  showstringspaces=false,                           % underline spaces within strings
  showtabs=false,                                   % show tabs within strings adding particular underscores
  stepnumber=1,                                     % the step between two line-numbers. If it's 1, each line will be numbered
  stringstyle=\color{mauve},                        % string literal style
  tabsize=2,                                        % sets default tabsize to 2 spaces
  title=\lstname                                    % show the filename of files included with \lstinputlisting; also try caption instead of title
}

\renewcommand\lstlistingname{Código} % muda o título padrão (Listings #) do pedaço de código para "Código #"


% codigo para que as referencias não sejam numeradas como um capitulo
\makeatletter
\renewcommand{\bibsection}{%
  \cleardoublepage
  \phantomsection
  \begin{center}
    \ABNTEXchapterfont
    \bfseries\fontsize{16pt}{24pt}\selectfont
    \MakeUppercase{\bibname}
  \end{center}%
  \vspace{2\baselineskip}%
  \addcontentsline{toc}{chapter}{\bibname}%
}
\makeatother


% ----------------------------------------------------------
% Alteração da forma de numeração de Tabelas e Figuras
% ----------------------------------------------------------

\counterwithout{figure}{section}
\counterwithout{table}{section}